\section{Variable definitions}\label{app:variables}

\begin{itemize}
    \item \textbf{firm cost share} (\eqref{eq:fcs}): pre-shock firm-level carbon cost share, time-invariant. Numerator: 2013--2016 average annual allowance shortfall in tCO$_2$ priced at the year-specific EUA. Denominator: 2010--2012 average total operating cost from Annual Accounts. Falls back to the earliest available three-year window for firms without 2010--2012 data.
    \item \textbf{nace exposure} (\eqref{eq:nace_exposure}): sales-weighted average firm cost share within a NACE 4-digit category, time-invariant; equals zero for categories with no ETS sellers.
    \item \textbf{share top}: annual share of the most-exposed ETS supplier $j^*(b,n)$ in cell $(b,n)$'s total spending in that NACE 4-digit category.
    \item \textbf{pair exposure}: $\text{firm cost share}_{j^*(b,n)} \times \bar{s}_{j^*(b,n)}$, where $\bar{s}_{j^*(b,n)} = \overline{\text{spending from } b \text{ to } j^*}_{2010..2014} / \overline{\text{total inputs of } b}_{2010..2014}$. Units: percentage of buyer's total input bill at risk if $j^*$'s carbon shock fully passes through and the buyer does not substitute.
    \item \textbf{Post}: $\mathbf{1}[t \geq 2015]$.
    \item \textbf{regulated}: $\mathbf{1}[\text{nace exposure}_n > 0]$.
\end{itemize}

\section{Pass-through: CMdG-style regulated-vs-unregulated event study}\label{app:passthrough_cmdj}

We complement the headline panel-LP elasticity in \S\ref{sec:passthrough} with an event-study replication of \citet[eq.~1]{CosterMejeanDiGiovanni2024} on Belgian monthly NACE 4-digit PPI data. The specification is exactly CMdG's:
\begin{equation}
    \log \text{PPI}_{s,t} \;=\; \sum_{\tau \neq 2004} \beta_\tau \,\mathbf{1}[s \in \text{regulated}]\,\mathbf{1}[\text{year}(t) = \tau] \;+\; \alpha_{s, \,\text{month}(t)} \;+\; \delta_{\text{year}(t)} \;+\; \varepsilon_{s,t}
    \label{eq:passthrough_cmdj}
\end{equation}
estimated on the 2000m1--2022m12 monthly panel with sector $\times$ calendar-month fixed effects $\alpha_{s, \text{month}(t)}$ (controlling for sector-specific seasonality, as in CMdG) and year fixed effects $\delta_{\text{year}(t)}$. Standard errors are clustered on NACE 4-digit. ``Regulated'' is CMdG's tight definition (seven NACE 2-digit codes \{17, 19, 20, 23, 24, 25, 35\}, per their Table~A.5 col~(1)/(4)). The reference year is 2004, the last pre-ETS year.

This specification differs from the headline panel-LP in \S\ref{sec:passthrough} in two ways. First, identification is binary (regulated vs unregulated NACE 2-digit) rather than continuous in $\omega$. Second, no carbon-policy shock is used: the variation comes from the relative time-series of regulated vs unregulated PPI, not from regulatory news events. The CMdG specification is thus less structural — it documents that prices in regulated sectors rose faster than in unregulated sectors over Phase~II and beyond, but it does not isolate the EU ETS as the cause. We report it because it is the standard reference in the literature \citep{CosterMejeanDiGiovanni2024}, and because it provides a sanity check on the panel-LP result at a different identifying margin.

\begin{figure}[H]
    \centering
    \includegraphics[width=0.85\textwidth]{phase1_figure1_cmdj_style.png}
    \caption{Belgian replication of \citet{CosterMejeanDiGiovanni2024} Figure~1: year-by-year coefficients $\widehat{\beta}_\tau$ from equation~(\ref{eq:passthrough_cmdj}). Sample: 2000m1--2022m12 Belgian NACE 4-digit monthly PPI panel; reference year 2004. Vertical bars are 95\% pointwise confidence intervals (clustered on NACE 4-digit). Shaded bands denote ETS Phases~1 and~3.}
    \label{fig:cmdj_ppi_eventstudy}
\end{figure}

\paragraph{Result.} The year-by-year coefficients show a clear structural break around 2005 (Phase~I start), with regulated PPI rising 2--4 log-points relative to unregulated by 2005--2007, then 7--10 log-points through Phase~II, and 9--14 log-points by Phase~IV. Phase-aggregated coefficients (relative to the 2000--2004 pre-ETS reference period):
\[
    \widehat{\beta}_{\text{Phase 1}} = +0.088^{***}, \quad
    \widehat{\beta}_{\text{Phase 2}} = +0.138^{***}, \quad
    \widehat{\beta}_{\text{Phase 3}} = +0.098^{**}, \quad
    \widehat{\beta}_{\text{Phase 4}} = +0.193^{***}.
\]
A pre-trend is visible: pre-ETS coefficients drift from $-0.077$ in 2000 toward zero at the 2004 reference, suggesting the regulated-vs-unregulated PPI gap was already closing before the EU~ETS began. Linear extrapolation of the pre-trend forward predicts $\widehat{\beta}_h \approx +0.08$ at $h = 2008$ from the trend alone; the actual Phase~II average ($+0.138$) is meaningfully above that extrapolation, but the pre-trend leaves the design less than a pristine parallel-trends test. The CMdG specification on Belgian data is therefore consistent with positive and economically large pass-through, but with a parallel-trends caveat that motivates the headline reliance on the Käenzig-shock identification of \S\ref{sec:passthrough} as the cleaner identifying source.

\section{Pass-through: LP-IV invertibility-robust robustness}\label{app:passthrough_lp_iv}

The headline pass-through specification (\eqref{eq:passthrough_strategy1}) uses Käenzig's VAR-identified Shock as the regressor in a panel local projection. Identification of the Shock from \citet{kanzig2023unequal}'s external-instruments SVAR rests on (partial) invertibility. We follow \citet[Appendix C.5]{kanzig2023unequal} in supplementing the headline with a local-projection instrumental-variables (LP-IV) exercise that does not require invertibility, using the raw refined Surprise series as the instrument:
\begin{equation}
    \log P_{s, m+h} - \log P_{s, m-1} \;=\; \theta_h \cdot \omega_{s, t(m)} \cdot \big(\log p^{\text{EUA}}_m - \log p^{\text{EUA}}_{m-1}\big) \;+\; \sum_{j=1}^{6} \phi_{h,j} \log P_{s, m-j} \;+\; \alpha_s \;+\; \delta_m \;+\; \xi_{s,m,h}
    \label{eq:passthrough_lp_iv}
\end{equation}
estimated by 2SLS, with $\omega_{s, t(m)} \cdot \text{CPSurprise}_m$ instrumenting for $\omega_{s, t(m)} \cdot \Delta \log p^{\text{EUA}}_m$. As prescribed by \citet{kanzig2023unequal}, the impulse horizon is restricted to $h \leq 12$ months because the LP-IV signal-to-noise ratio at longer horizons is too low to credibly identify the carbon-policy effect against the myriad other shocks affecting EUA prices. Inference uses cluster-robust standard errors at NACE 4-digit, with lag-augmentation as in \citet{kanzig2023unequal}'s appendix specification.

\begin{figure}[H]
    \centering
    \includegraphics[width=0.85\textwidth]{phase3_passthrough_lp_iv_robust.pdf}
    \caption{LP-IV invertibility-robust impulse response $\widehat{\theta}_h$ from equation~(\ref{eq:passthrough_lp_iv}). $\omega_{s, t(m)} \cdot \text{CPSurprise}_m$ instruments for $\omega_{s, t(m)} \cdot \Delta \log p^{\text{EUA}}_m$. Sample: 2010m1--2019m12 Belgian NACE 4-digit monthly panel, 134 sectors. Solid line: point estimate; shaded band: 95\% confidence interval (cluster-robust at NACE 4-digit). Lag-augmentation with six lags of $\log P_{s, m-j}$. Horizon capped at 12 months per \citet[App.~C.5]{kanzig2023unequal}.}
    \label{fig:passthrough_lp_iv}
\end{figure}

\paragraph{Reading.} Confidence intervals are wide enough to include zero at every horizon, which matches the pattern \citet{kanzig2023unequal} reports for the LP-IV variant of his own HICP IRF (his Figure C.20). The point estimates do not contradict the headline panel-LP IRF — the headline's $\widehat{\gamma}_h$ at horizons up to $h = 12$ lies within the 95\% LP-IV band — but the LP-IV is too underpowered on the Belgian panel to deliver an independent point estimate. As Käenzig notes, this is a known limitation of LP-IV with high-frequency policy surprises against macroeconomic outcomes: the surprise series captures only the regulatory-news component of EUA-price variation, and at the monthly frequency this component is small relative to the cumulative effect of other (gas-market, demand, weather) shocks driving EUA. The robustness exhibit confirms that the headline conclusion does not rest on the SVAR's invertibility assumption in a way that LP-IV would reject.
